\documentclass[11pt, oneside]{article} 
\usepackage{geometry}
\geometry{letterpaper} 
\usepackage{graphicx}
	
\usepackage{amssymb}
\usepackage{amsmath}
\usepackage{parskip}
\usepackage{color}
\usepackage{hyperref}

\graphicspath{{/Users/telliott/Github-Math/figures/}}
% \begin{center} \includegraphics [scale=0.4] {gauss3.png} \end{center}

\title{Cycloid}
\date{}

\begin{document}
\maketitle
\Large

%[my-super-duper-separator]

Imagine a bicycle with one tire marked at a particular point on the rim, say with fluorescent paint or a small light.  Time starts at $t = 0$ with that point $P$ in contact with the $x$ axis at $(0,0)$.  Then the bike rolls to our right.  As the tire rotates the fixed point $P$ on the rim traces a curve
\begin{center} \includegraphics [scale=0.6] {cycloid.png} \end{center}

We want to find equations that give the point $P$ as a function of time.  We will parametrize the curve, yielding parametric equations $x(t)$, $y(t)$.

The second diagram  shows the angle through which the wheel has turned as $\theta$, but we will use $t$ for $\theta$ here.  

\begin{center} \includegraphics [scale=0.5] {cycloid2.png} \end{center}
The $x$ displacement of the vertical straight down from the center of the tire is just $at$, where $a$ is the radius of the wheel, it is equal to the arc on the circumference of the wheel from the point which is currently in contact with the ground, going around up to $P$.

It is reasonably easy to derive the desired parametric equations, using vectors, especially once you know the answer.  For $x$, we have the vector that goes from $(0,0)$ to the contact point with the ground.  As indicated in the figure, that is $at$.  

We need to subtract the distance $a \ sin\ t$ from that.  Basically the rationale is that the motion is a standard parametric circle which has been rotated by $90$ degrees clockwise and then inverted.  The rotation changes cosine to sine, and the inversion brings the subtraction.

It's easier to see for $t < \pi/2$, but it is true always.  Check some other values of $t$ like $\pi$ or $3\pi/2$ to confirm.  This is the usual circular motion.

For $y$, we have a constant factor of $a$ above the $x$ axis, then the additional displacement is $-a \ cos \ t$.  So for $t=0$ we have the additional displacement is -a (we were on the ground), for $t=\pi/2$ it is zero, and for $t=\pi$ it is plus $a$ for a total of $2a$.

The parametric equations are then
\[ x(t) = at - a \sin t \]
\[ y(t) = a - a \cos t \]
Taking derivatives:
\[ x'(t) = a - a \cos t \]
\[ y'(t) = a \sin t  \]
We can get $dy/dx$, what we usually call $y'$, by simple division:
\[ y' = \frac{dy}{dx} = \frac{a \sin t }{a - a \cos t} \]
\[ = \frac{\sin t}{1 - \cos t} \]

Simmons uses half-angle formula like so:
\[ \sin t = 2 \sin t/2 \cos t/2 \]
\[ \cos t = \cos^2 t/2 - \sin^2 t/2 \]
\[ = 1 - 2 \sin^2 t/2 \]
So the ratio is
\[ y' = \frac{ 2 \sin t/2 \cos t/2}{2 \sin^2 t/2} = \cot t/2 \]

\subsection*{aside about Archimedes}

It struck me that $\cot t/2$ is one of the terms in the relationship Archimedes used in approximating $\pi$.  There we had
\[ \cot 2 \theta + \csc 2 \theta = \cot \theta \]
\[ \cot t + \csc t = \cot t/2 \]
So somehow, if what we have above is correct
\[ \frac{\sin t}{1 - \cos t} = \cot t/2 = \cot t + \csc t \]
Factor the right-hand side
\[  = \frac{1}{\sin t} (\cos t + 1) \]
Then
\[ \sin^2 t = (1 - \cos t)(\cos t + 1) \]
\[ = 1 - \cos^2 t \]
and one more step gives our favorite identity.

\subsection*{Area under the arc}
We want
\[ A = \int_{t=0}^{t=2\pi} y \ dx \]
\[ = \int_{t=0}^{t=2\pi} (a - a \cos\ t) (a - a \cos\ t) \ dt \]
\[ a^2\int_{t=0}^{t=2\pi} (1 - \cos\ t) (1 - \cos\ t) \ dt \]
\[ a^2\int_{t=0}^{t=2\pi} (1 - 2 \cos\ t + \cos^2\ t) \ dt \]
If you don't remember the result for $\int \cos^2 t \ dt$, we will derive the double angle formula later and convert from $\sin^2$ to $\cos^2$.  Write:
\[ A = a^2 ( t - 2 \sin \ t + \frac{1}{2}t + \frac{1}{4} \sin 2t ) \ \bigg|_0^{2\pi} \]
\[ a^2 ( 2\pi - 0 + \pi + 0 - 0 + 0 - 0 - 0    ) = 3\pi a^2 \]
A very simple answer.

\subsection*{Roberval}
The cycloid curve was not known to the ancients, as far as our sources have anything to say.  The first mention is around 1500, and by the time of Galileo a century later, it was something he thought about quite a bit.  Galileo is said to have used Archimedes method of cutting out shapes and weighing them) to show that the area was approximately $3$ times that of the generating circle.

The cycloid became very popular in the decades before Newton.  Roberval came up with a clever application of Cavlieri's principle, to find the area under the curve.

The idea is to draw a second curve, called the \emph{companion curve}, by adding a half-circle to the $x$-values on the cycloid curve, as shown.

The areas marked by horizontal lines are equal, by Cavalieri's principle.
\begin{center} \includegraphics [scale=0.2] {cycloid_comp.png} \end{center}

\url{https://maa.org/sites/default/files/pdf/cmj_ftp/CMJ/January%202010/3%20Articles/3%20Martin/08-170.pdf}

But the curve $AQD$ apparently divides the rectangular area $ACDB$ in half, by symmetry.  (It is actually the sine curve from $-\pi/2$ to $\pi/2$, just scaled up by a factor of $a$ in both directions).  Since the rectangle has width $\pi a$ and height $2a$, its area is $2 \pi a^2$, so one-half that is $\pi a^2$.

The area under the cycloid curve is then three-quarters of the total, which (between $0 \rightarrow \pi$) is $(3/2) \pi a^2$, and the area under one complete lobe is twice that or $3\pi a^2$.

Let's see if we can figure out a parametric equation for the companion curve.  $x(t)$ for the cycloid was $x(t) = at - a \sin t$.  The companion curve gets an additional length in the $x$-direction of $a \sin t$.  So $x(t) = at$.  That was easy!

$y(t)$ is unchanged.  Our integral is just $\int y \ dx$ and $dx = a \ dt$ so
\[ A = \int (a - a\cos t) \ a \ dt \]
The bounds are $0 \rightarrow \pi$ so finally
\[ A = a^2 \int_0^{\pi}  1 - \cos t \ dt \]
\[ = a^2 \ [ \ t - \sin t \ ] \bigg |_0^{\pi} \]
\[ = \pi a^2 \]
which matches what we had before.

\subsection*{slope}
Descartes gave the slope of the tangent line to the cycloid by the following construction:
\begin{center} \includegraphics [scale=0.2] {cycloid_slope.png} \end{center}
Given $P$ on the cycloid, draw $PE$ horizontally to find $E$ on the circle, then draw $PQ$ parallel to $EC$.  The tangent is perpendicular to $PQ$.

I like this construction a lot because it reminds me of what Roberval did with the parabola.  For any $t$, the distance moved horizontally is equal to the distance along the rim of the generating circle.  The tangent should add to or subtract from both values at the same rate.  

In other words, I expect that the tangent should be the average of the tangent to the circle at $E$ and the horizontal, namely one-half of the former.

If $C$ is taken to be $(0,0)$, the coordinates of $E$ are $(x',y)$, where
\[ x' = -a \sin t  \]
and the slope of $PQ \parallel EC$ will be 
\[ \frac{0 - (a - a \cos t)}{0 - (-a \sin t)} = \frac{\cos t - 1}{\sin t} = \cot t - \csc t \]
The slope of the tangent is the negative inverse
\[ \frac{\sin t}{1 - \cos t} \]
which matches.

Recall that we had before that
\[ y' =  \cot t/2 \]
\begin{center} \includegraphics [scale=0.2] {cycloid_slope.png} \end{center}
which seems like a problem at first but $t \ne \angle ECQ$.  

If we had drawn the center of the circle as $O$, then $t$ would be $\angle EOC$, and since $\triangle EOC$ is isosceles
\[ t + 2 \cdot \angle OCE = 180 \]
\[ t/2 + \angle OCE = 90 \]
but $\angle OCE$ ($\angle DCE$) and $\angle ECQ$ are also complementary so
\[ t/2 = \angle ECQ = \angle EPQ \]

The angle made by the tangent ($\angle HPE$) is complementary, so its tangent is the cotangent of $t/2$.

Note also that $\angle EDC = \angle ECQ = \angle EPQ$, so it is also complementary to the tangent angle.

\subsection*{Arc length}
We wish to determine the arc length and area under the curve for one complete revolution of the wheel.

We want to use a slightly different version of the usual formula for arc length

\[ ds^2 = dx^2 + dy^2 \]
\[ (\frac{ds}{dt})^2 = (\frac{dx}{dt})^2 + (\frac{dy}{dt})^2 \]
\[ ds = \sqrt{(\frac{dx}{dt})^2 + (\frac{dy}{dt})^2} \ dt \]
\[ = \sqrt{(a - a \cos\ t)^2 + (a \sin\ t)^2} \ dt  \]
This expands to
\[ a \sqrt{1 - 2 \cos\ t + \cos^2t + \sin^2t } \ dt \]
\[ =  a \sqrt{2 - 2 \cos \ t} \ dt\] 
The length is
\[ L = \int_0^{2\pi} a \sqrt{2 - 2 \cos \ t} \ dt\]
\[ = a \sqrt{2} \ \int_0^{2\pi} \sqrt{1 - \cos \ t} \ dt\]

\subsection*{double angle}

\[ \cos (s-t) = \cos s \cos t + \sin s \sin \ t \]
(check:  if $s=t$ then $\cos 0 = 1$, which is correct).

So
\[ \cos (s+t) = \cos s \cos t - \sin s \sin \ t \]
Let $s = t$ and $u = 2s$, then
\[ \cos 2s = \cos u = \cos^2 \ (\frac{u}{2}) - \sin^2 \ (\frac{u}{2}) \]
\[ \cos u = 1 - \sin^2 \ (\frac{u}{2}) - \sin^2 \ (\frac{u}{2}) \]
\[ 2 \sin^2 \ (\frac{u}{2}) = 1 - \cos u \]
$u$ is just a dummy variable, so we can switch back to $t$
\[ 2 \sin^2 \ (\frac{t}{2}) = 1 - \cos t \]

\subsection*{finishing up}
We have that 
\[ L = a \sqrt{2} \ \int_0^{2\pi} \sqrt{1 - \cos \ t} \ dt\]
And
\[ 1 - \cos t = 2 \sin^2(\frac{t}{2}) \]
\[ \sqrt{1 - \cos t} = \sqrt{2} \sin(\frac{t}{2}) \]
So
\[ L = a \sqrt{2} \ \int_0^{2\pi} \sqrt{2} \sin \ (\frac{t}{2}) \ dt\]
\[  = 2a  \ \int_0^{2\pi} \sin \ (\frac{t}{2}) \ dt\]
\[ = 2a \ (-2) \cos \ (\frac{t}{2})\ \bigg |_0^{2\pi} \]
\[ = -4a \ (\cos \ \pi - \cos \ 0) \]
\[ = -4a \ (-1 - 1) = 8a\]
Also a very simple answer to the problem.

\subsection*{tautochrone}

\url{https://en.wikipedia.org/wiki/Tautochrone_curve}

According to Wolfram

\url{https://mathworld.wolfram.com/TautochroneProblem.html}

``The problem of finding the curve down which a bead placed anywhere will fall to the bottom in the same amount of time.''  It turns out that the cycloid satisfies this condition, if it is placed upside down.

The derivation in Wolfram is complicated by having to keep track of the difference between the parameters $t$ and $\theta$.  

Up till now we have considered them to be the same, with time synchronized to the angle, so $t$ is just our symbol for the angle, replacing $\theta$, but now we will have arbitrary $t$ with $0 \le t \le 1$ and $0 \le \theta \le \pi$.

The parametric equations for the cycloid in terms of the angle $\theta$:
\[ x = a (\theta - \sin \theta) \]
\[ \frac{dx}{d \theta} = a(1 - \cos \theta) \]
\[ y = a (1 - \cos \theta) \]
\[ \frac{dy}{d \theta} = - a \sin \theta) \]
So
\[ x'(t) = \frac{dx}{dt} = \frac{dx}{d \theta} \frac{d \theta}{dt} \]
\[ = a(1 - \cos \theta) \ \frac{d \theta}{dt} \]
\[ y'(t) = \frac{dy}{dt} = \frac{d}{d \theta} \frac{d \theta}{dt}  \]
\[ = -a \sin \theta \ \frac{d \theta}{dt} \]
From which we easily obtain $dx$ and $dy$ and then the arc length is
\[ ds = \sqrt{dx^2 + dy^2} \]
\[ =  \sqrt{a^2(1 - \cos \theta)^2 + a^2 \sin^2 \theta}  \ d \theta \]
(as before)
\[ = \sqrt{2} a (1 - \cos \theta) \ d \theta \]

The actual derivation just follows Abel's use of conservation of energy:
\[ \frac{1}{2} mv^2 = mgy \]
(note that $y$ is positive up and $v$ is positive down).
\[ v = \frac{ds}{dt} = \sqrt{2gy} \]
Combining these results
\[ dt = \frac{ds}{\sqrt{2gy}}  \]
\[ = \frac{ \sqrt{2} a (1 - \cos \theta) \ d \theta}{\sqrt{2g a (1 - \cos \theta)}} \]
\[ = \sqrt{\frac{a}{g}} \ d \theta \]
\[ t = \int_0^1 \ dt = \sqrt{\frac{a}{g}} \  \int_0^{\pi} d \theta \]
\[ t = \pi \  \sqrt{\frac{a}{g}} \]

I'll just refer you to the Wolfram article to see that they get the same answer starting from
\[ v = \frac{ds}{dt} = \sqrt{2g(y - y_0)}  \]


\end{document}